%!TEX Root = ../main.tex
\chapter{Localization}
\label{chap:localization}
In this chapter, we will discuss one of the most
important construction in commutative algebra: localization.
The motivation of localization is the construction of
the fraction field $\Frac A$ from an integral domain $A$.
We would like to generalize this construction
to specify which element can and cannot be
the denominator.
Soon, we will quickly see that this construction
allows us to ``zoom into'' the ring $A$ and study them
at the level of either open sets (read: $D(f)$
for $f\in A$) or points (read: prime ideals $\pp$).
In the latter case, we get a local ring,
which is usually easier to study.

\section{Localization of Rings}
\subsection{Definition and Explicit Construction}
We will define localization by
its universal property and then construct it later.
\begin{definition}[Multiplicative set]
    Given a ring $A$, a subset $S\subseteq A$ is 
    a \vocab{multiplicative set} 
    if and only any finite product of elements 
    in $S$ is in $S$.
    In particular, the empty product, $1$, must be in $S$.
\end{definition}
\begin{definition}[Localization]
    \notation{$S^{-1}A$}{localization of a ring}
    Let $A$ be a ring.
    Let $S\subset A$ be a multiplicative set. Then,
    the \vocab{localization of ring} $S^{-1}A$ is a ring with
    map $\phi : A\to S^{-1}A$ such that
    \begin{itemize}
	\item $\phi(s)$ is a unit for all $s\in S$.
	\item \textbf{Universal property.}
        For any $A$-Algebra $B$ with map
	$\psi : A\to B$ such that $\psi(s)$ is a unit
	for all $s\in S$, the map $\psi$ factors through
	$S^{-1}A$: there is a unique map $
	\widetilde\psi : S^{-1}A\to B$
	such that $\psi = \widetilde\psi \circ\phi$.
    \begin{center}
    \begin{tikzcd}
        S \arrow[rr] & & B^\times \\[-8pt]
        A \arrow[rr,"\psi"] \arrow[rd,"\phi",swap] & & B \\
        & S^{-1}A \arrow[ru, "\widetilde\psi",swap]
    \end{tikzcd}
    \end{center}
    \end{itemize}
\end{definition}
\begin{theorem}
    $S^{-1}A$ exists and is unique.
\end{theorem}
As usual with universal property, the uniqueness
is automatic.

To construct localization, let the set of fractions be
$$A\times S = \left\{\frac as := (a,s)\right\}.$$
Of course, we have to handle something like
$\frac 12 = \frac 36$. We thus define an equivalence relation
$\frac as \sim \frac bt$ if and only if $at-bs=0$.
However, this is not an
equivalence relation (when $A$ is not a domain).

To salvage this, we make a slightly different construction:
We say that
$$\frac as\sim\frac bt\iff u(at-bs)=0\text{ for some } u\in S,$$
and one can check that this is an equivalence relation,
and we may let
$$S^{-1}A = \frac{S\times A}{\sim}.$$

We may define addition and multiplication by
the usual high school fraction operations:
\begin{align*}
    \frac as + \frac bt &= \frac{at+bs}{st} \\
    \frac as \cdot\frac bt &= \frac{ab}{st}.
\end{align*}
When writing this, one needs to check that the 
equivalence relation behaves
well with respect to addition and multiplication.
This is very tedious and hence omitted.
We can also define the inclusion $A\to B$ as
$\phi(a)=\frac a1$.
\begin{proposition}
    The construction of $S^{-1}A$ above
    satisfies universal property.
\end{proposition}
\begin{proof}
    Given $\psi : A\to B$ such that $\psi(s)^{-1}$ exists
    for all $s\in S$,
    one can define the map
    \begin{align*}
	\widetilde\psi : S^{-1}A &\to B \\
	\frac as &\mapsto \psi(a)\psi(s)^{-1}.
    \end{align*}
    One can check that this map respects the 
    equivalence relation: 
    if $u(at-bs)=0$, then $\psi(a)\psi(s)^{-1}
    = \psi(b)\psi(t)^{-1}$. To do so, note that
    $$\psi(u)\Big(\psi(a)\psi(t) - \psi(b)\psi(s)\Big) = 0,$$
    and since $\psi(u)$, $\psi(t)$, $\psi(s)$ are units,
    one can multiply by $\psi(u)^{-1}\psi(s)^{-1}\psi(t)^{-1}$
    on both sides to get that $\widetilde\psi$ is well-defined.

    Finally, it's obvious that $\widetilde\psi$
    satisfies the universal property.
\end{proof}
%
%
%
%
%
\subsection{Examples}
\begin{example}[Fraction Field]
    \notation{$\Frac A$}{fraction field}
    If $A$ is a domain and $S=A\setminus\{0\}$, then
    $$S^{-1}A  =: \Frac A$$
    is a field.

    For example, if $A=\ZZ$ and $S=\ZZ\setminus\{0\}$,
    then $\Frac \ZZ = S^{-1}\ZZ = \QQ$.
\end{example}
\begin{example}[Localization at one element]
    Take $A=\mathbb C[x]$ and $S=\{(x-1)^k : k\geq 0\}$.
    Then,
    $$S^{-1}\mathbb C[x]
    = \left\{\frac{g(x)}{(x-1)^k} : g\in\CC[x]\right\}.$$
    Geometrically, this ring is a ring of function defined everywhere at $\CC\setminus\{1\}$,
    thus allowed to have poles at $1$.
\end{example}
% Let's make a similar example but with two variables.
% \begin{example}
%     Take $A=\mathbb C[x,y]$ and $S=(y-x^2)^k$.
%     Then, $S^{-1}\mathbb C[x,y]$ is a set of
%     two-variable rational function defined at
%     $\mathbb C^2 \setminus (\text{parabola }y=x^2)$.
% \end{example}
\begin{example}[Two common localizations, in $\ZZ$]
    Take $A=\mathbb Z$. There are two common localization
    about some prime, say $5$.
    \begin{itemize}
	\item Take $S=\{5^k : k\geq 0\}$, we get
	    $\mathbb Z\left[\frac 15\right]\in\mathbb Q$.
	\item Take $S=\{n : 5\nmid n\}$ to get
	    $S^{-1}\mathbb Z = \mathbb Z_{(5)}
	    = \left\{\frac ab : 5\nmid b\right\}$.
    \end{itemize}
\end{example}
\begin{example}[The two most important examples of localization]
    Let $A$ be a ring. Then, there are two important examples.
    \begin{alphenum}
	\item \vocab[localization away from $f$]{Localization away from $f$.}
        \notation{$A[f^{-1}]$, $A_f$}{localization
        away from $f$}
        For any $f\in A$, take
	    $S = \{1,f,f^2,\dots,\}$ to get
	    $$S^{-1}A = A_f = A[f^{-1}]
	    = \frac{A[x]}{(xf-1)}.$$
	\item \vocab[localization at prime ideal $\pp$]{Localization at prime ideal $\pp$.}
        For any prime ideal $\mathfrak p$, take
	    $S=A\setminus\mathfrak p$ to get
	    $$S^{-1}A = A_{\mathfrak p}.$$
        \notation{$A_\pp$}{localization at $\pp$}
	    When $A$ is a domain and $\mathfrak p=(0)$,
	    we get the fraction field.
    \end{alphenum}
\end{example}
%
%
%
%
\subsection{Ideals of Localization}
% First, let's state an obvious fact.
% \begin{proposition}
%     If $A$ is a domain, then
%     $A\subset S^{-1}A\subset\Frac A$.
% \end{proposition}
% Now, let's look at how localization works in quotient.
% \begin{proposition}
%     Let $A$ be a ring, $S$ be multiplicatively closed,
%     
% \end{proposition}
% \begin{proof}
%     Homework 5.
% \end{proof}
% Here is an example of applying the above proposition.
% \begin{example}
%     Let $A=\mathbb C[x,y]/(xy)$ and 
%     $S=\{1,x,x^2,\dots\}$. Then,
%     $\overline S=\{1,\overline x,\overline x^2,
%     \dots\}$, so
%     $$\overline S^{-1}A = \frac{\mathbb C[x,y]\left[\frac 1x\right]}
%     {(xy)} = \frac{\mathbb C\left[x,y,\frac 1x\right]}
%     {(xy)}.$$
% \end{example}
Let's discuss how ideals behave under localization.
From the map
$\phi : A\to S^{-1}A$, the ideal $I$ extends to ideal
\notation{$S^{-1}I$}{localization of ideal $I$}
$$S^{-1}I := \left\{\frac is : i\in I\right\} 
\ \subseteq\ S^{-1}A.$$

It turns out that every ideal of $S^{-1}A$ is of this form.
\begin{proposition}[Every ideal is extended ideal]
    \label{prop:ideal_is_extended}
    For all ideal $J\subseteq S^{-1}A$,
    we have $S^{-1}(\phi^{-1}(J)) = J$.
    
    \noindent In particular, every ideal in $S^{-1}A$ is of the form 
    $S^{-1}I$ for some ideal $I\subseteq A$.
\end{proposition}
\begin{proof}
For all $a\in A$ and $s\in S$,
each of the following is equivalent to the next.
\begin{itemize}
\item $\frac as\in J$
\item $\frac a1\in J$. (Proof: $s$ is a unit in $S^{-1}A$).
\item $a\in\phi^{-1}(J)$.
\item $\frac as\in S^{-1}(\phi^{-1}(J))$. \qedhere
\end{itemize}
\end{proof}

\begin{theorem}[Prime ideals of localization]
\label{thm:prime_ideal_localize}
    Let $A$ be a ring, $S$ multiplicatively closed,
    and $\phi : A\to S^{-1}A$ be the canonical map.
    There is a correspondence between
    \begin{align*}
	\{\text{prime ideals of } S^{-1}A\} &\leftrightarrow 
	    \{\text{prime ideal } \mathfrak p\subset A
	    \text{ such that } \mathfrak p\cap S = \emptyset\}
	\\
	\mathfrak q &\mapsto \phi^{-1}\frakq \\
	S^{-1}\mathfrak p &\mapsfrom \mathfrak p.
    \end{align*}
\end{theorem}
\begin{proof}
    The proof is in several steps.
    \begin{itemize}
    \item \textbf{The first map is well-defined.}
    $\phi^{-1}\frakq$ is prime because
    contraction of prime ideal is prime.
    Moreover, $\phi^{-1}\frakq$ does not meet $S$ because
    if $s\in S$ and $s\in \phi^{-1}\frakq$,
    then $\phi(s)\in\frakq$, but $\phi(s)=\frac s1$ is a unit
    in $S^{-1}A$, so $\qq=S^{-1}A$, a contradiction.
    
    \item \textbf{The second map is well-defined.}
    Let $\pp$ be a prime ideal that does not meet $S$.
    We first show that for all $a\in A$ and $s\in S$,
    \begin{equation}
        \label{eq:localize_ideal_condition}
        \tfrac as \in S^{-1}\pp\ \iff\ a\in\pp.
    \end{equation}
    The direction $(\Leftarrow)$ is clear.
    For the direction $(\Rightarrow)$, suppose that 
    $\tfrac as = \tfrac bt$ for some $b\in\pp$ and $t\in S$.
    Then, there exists $u\in S$ such that 
    $u(at-bs) = 0$. Reducing modulo $\pp$ gives 
    $uat\in\pp$. However, $u,t\in S$, so they are not in $\pp$,
    so $a\in \pp$.

    Finally, if $\tfrac as$ and $\tfrac bt$ are in 
    $S^{-1}A$ such that $\tfrac as \cdot\tfrac bt\in S^{-1}\pp$,
    then by \eqref{eq:localize_ideal_condition}, 
    $ab\in \pp$, and so either $a\in\pp$ or $b\in\pp$.
    Thus, either $\tfrac as\in S^{-1}\pp$ 
    or $\tfrac bt\in S^{-1}\pp$. Hence, $S^{-1}\pp$ is prime.
    \item \textbf{For all prime $\qq\subseteq S^{-1}A$, 
    we have $S^{-1}(\phi^{-1}\qq)=\qq$.}
    Follows from \Cref{prop:ideal_is_extended}.
    \item \textbf{For all prime $\pp\subseteq A$
    that does not meet $S$, we have $\phi^{-1}(S^{-1}\pp)
    = \pp$.}
    Indeed, each of the  following is equivalent to the next
    \begin{itemize}
	\item $a\in\phi^{-1}(S^{-1}\mathfrak p)$
	\item $\frac a1\in S^{-1}\mathfrak p$
	\item $a\in \mathfrak p$. (by \eqref{eq:localize_ideal_condition}) \qedhere
    \end{itemize}
    \end{itemize}
\end{proof}
Recall that spectrum of a ring is the topological space
of prime ideals of a ring.
\begin{corollary}[Spectrum of localization
    away from an element]
    For all $f\in A$, we have
    $$\Spec A\left[1/f\right]
    \simeq D(f) = \{\mathfrak p : f\notin\mathfrak p\},$$
    which is an open subset of $\Spec A$.
\end{corollary}
\begin{corollary}[Spectrum of localization
    at a prime ideal]
    For all prime ideal $\mathfrak p$,
    we have 
    \begin{alphenum}
    \item $\Spec A_{\mathfrak p}
    \simeq \{\mathfrak q : \mathfrak q \subset\mathfrak p\}$.
    \item $A_\pp$ is a local ring with unique maximal ideal
    $\pp A_\pp$.
    \end{alphenum}
\end{corollary}
We now explain where the term localization comes from.
\begin{itemize}
    \item $\Spec A[1/f]$ is the open subset $D(f)$ of $\Spec A$ of points where $\tfrac 1f$ is defined. 
    In other words, functions in $A[1/f]$ are rational functions
    defined on the open set $D(f)$.
    We are localizing \textbf{away from}
    points that $f$ vanishes, and this explains
    why it gets the name ``localization away from $f$''.
    \item In contrast, $\Spec A/(f)$ is $V(f)$,
    which is complement of $D(f)$.
    \item $\Spec A_\pp$ is the intersection 
    of all open sets containing $\pp$.
    Thus, functions (read, elements) in $A_\pp$ 
    corresponds to rational functions 
    defined on some neighborhood of $\pp$.
    In particular, we are essentially zooming into 
    just point $\pp$, and this is why it gets the name
    ``localization at $\pp$''.
\end{itemize}
To drive the point home, let us draw picture of 
$\Spec \CC[x]_x = \Spec \CC[x,x^{-1}]$
and $\Spec \CC[x]_{(x)}$.
\begin{itemize}
\item All prime ideals of $\Spec \CC[x,x^{-1}]$
are $(0)$ and $(x-\alpha)$ for all $\alpha\in\CC\setminus\{0\}$.
Thus, it is an affine line minus point $0$.
\item All prime ideals of $\Spec \CC[x]_{(x)}$
are $(0)$ and $(x)$. 
Thus, $\Spec\CC[x]_{(x)}$ has only two points,
and all closed sets are $\emptyset, \{(x)\}, \{(0), (x)\}$.
\end{itemize}
\begin{center}
\begin{asy}
size(7cm,0);
pair A = (-8,0); pair B = (8,0);
draw(A--B, blue, Arrows);
dot("$(x-3)$", (3,0), dir(-90), blue);
dot("$(x+2)$", (-2,0), dir(-90), blue);
dot("$(0)$", (0,1), dir(90), red);
filldraw(circle((0,0),0.25), white, blue+1.5);
draw( (-6,0)..(-3,0.5)..(-1,1)--(1,1)--(3,0.5)--(6,0), dotted+red);
label("$\operatorname{Spec} \mathbb C[x,x^{-1}]$", A, dir(-90), blue);
\end{asy}
\quad 
\begin{asy}
size(2cm);
draw( (0,0)--(0,3), heavygreen);
label("$\operatorname{Spec} \mathbb C[x]_{(x)}$", (0,3), dir(90), blue);
draw( (-1,1)--(1,1), dotted+red );
dot("$(0)$", (0.2,1), dir(45), red);
dot("$(x)$", (0,0), dir(-90), blue);
\end{asy}
\end{center}

\subsection{Localization Commutes with Everything}
Localization commutes with everything you can ever imagine.
In this subsection, we state various theorems of this sort.
We omit all proofs because they
are all routine. (In fact, you should be able to prove
it yourself.)

As as setup, let $A$ be a ring, $S$ be a multiplicatively closed subset, and $\phi : A\to S^{-1}A$
be the canonical map $a\mapsto \tfrac a1$.
\begin{proposition}[Localization Commutes with Quotient]
Let $I$ be an ideal of $A$ and 
$\overline S$ be the image of $S$ in $A/I$, then
$$\overline S^{-1}(A/I) = \frac{S^{-1}A}{(\phi(I))}.$$
\end{proposition}
\begin{proposition}[Localization in Stages]
Let $T$ be another multiplicatively closed subset of $A$,
and let $U$ be the image of $T$ in $S^{-1}A$.
Then,
$$U^{-1}(S^{-1}A) = (ST)^{-1}A,$$
where $ST=\{st : s\in S, t\in T\}$.
\end{proposition}
\begin{proposition}[Localization Commutes with Nilradical]
    \label{prop:localization_nilradical}
    We have $S^{-1}(\Nil A) = \Nil(S^{-1}A)$.
    In particular, if $A$ is reduced
    (no nilpotent element), then so is $S^{-1}A$.
\end{proposition}
\begin{proposition}[Properties Preserved by Localization]
    \begin{alphenum}
    \item If $A$ is a domain, then $S^{-1}A$ is a domain.
    \item If $A$ is a PID, then $S^{-1}A$ is a PID.
    \item If $A$ is a UFD, then $S^{-1}A$ is a UFD.
    \end{alphenum}
\end{proposition}
\begin{exercise}
If you wish, prove all of the above theorems.
\end{exercise}
%
%
%
%
%
%
%
%
%
\section{Localization of Modules}
\subsection{Definition and Basic Properties}
The easiest way to define localization of modules
is to base change from $A$ to $S^{-1}A$.
\begin{definition}[Localization of modules]
    \notation{$S^{-1}M$}{localization of a module}
    Let $A$ be a ring and $S$ be multiplicatively closed.
    Then, $S^{-1}A$ is an $A$-algebra.
    Thus, one can define \vocab{localization of module} $M$
    as
    $$S^{-1}M := S^{-1}A\otimes_A M.$$
\end{definition}
This can be constructed explicitly as follows.
\begin{proposition}[Explicit construction]
    We have
    $$S^{-1}M = \frac{M\times S}{\sim},$$
    where the equivalence relation $\sim$ is defined by 
    $$\frac ms \sim\frac{m'}t\ \iff\ \text{ there exists } 
    u\in S\text{ such that } u(tm-sm')=0.$$
\end{proposition}
\begin{proof}
    There is nothing difficult in this proof.

    Let $M' = (M\times S)/\sim$ be the construction above.
    The map
    \begin{align*}
        S^{-1}A \times M &\to M' \\
        \frac as \times m &\mapsto \frac{am}s
    \end{align*}
    is bilinear, and hence descend down to the map
    $S^{-1}A\otimes_A M \to M'$. Clearly, this map is surjective.
    Now, we show that it is injective.
    First, we claim that every element in $S^{-1}A\otimes_A M$
    can be written in form $\tfrac 1s\otimes m$
    where $s\in S$ and $m\in M$. To prove this, suppose 
    we have one such element $x=\sum_{i=1}^n \tfrac{a_i}{s_i}
    \otimes m_i$. Let $s = s_1\cdots s_n$,
    $t_i = \prod_{j\neq i} s_j$. Then, we can rewrite
    $$x=\sum_{i=1}^n \frac{a_i}{s_i}\otimes m_i
    = \sum_{i=1}^n \frac{a_it_i}s \otimes m_i 
    = \sum_{i=1}^n \frac 1s \otimes a_it_im_i 
    = \frac 1s \otimes \left(\sum_{i=1}^n a_it_im_i\right)
    = \frac 1s\otimes m,$$
    as claimed.

    Now, suppose that $\frac 1s\otimes m$ is mapped to $0$.
    Then, $\tfrac ms = 0$ in $M'$, so there exists $u\in S$
    such that $um=0$. Therefore,
    $$\frac 1s \otimes m = u\left(\frac 1{us}\otimes m\right)
    = \frac 1{us}\otimes um = 0,$$
    implying that the map is injective.
\end{proof}
\subsection{Exactness}
Having localized modules, one can think of localization
as a functor between $A$-modules.
This functor is automatically right-exact
by \Cref{thm:tensor_is_right_exact}
because it is tensoring with $S^{-1}A$.
However, more is true.
\begin{theorem}[Localization is exact]
    \label{thm:localization_exact}
    The functor $S^{-1}\bullet$
    defined by $M\mapsto S^{-1}M$ is exact.
    In other words, an exact sequence 
    $0\to L\to M\to N\to 0$ gives the exact sequence
    $$0\longto S^{-1}L\longto S^{-1}M\longto S^{-1}N\longto 0.$$
\end{theorem}
\begin{proof}
    From \Cref{thm:tensor_is_right_exact},
    it suffices to show that $S^{-1}\bullet$
    preserves injectivity.
    Suppose $f : L\to M$ is injective.
    We need to show that $S^{-1}L\to S^{-1}M$ is injective.
    Suppose that there exist $l\in L$ and $s\in S$
    such that $\frac{f(l)}{s} = 0$.
    This implies that there exists $t\in S$ such that
    $tf(l)=0\in M$. 
    However, this implies that $f(tl)=0$, so $tl=0$,
    or $\frac ls=0$, so we are done.
\end{proof}
\begin{corollary}[Localization is flat]
    \label{cor:localization_flat}
    For any ring $A$ and any multiplicative set $S$,
    we have $S^{-1}A$ is a flat $A$-module.
\end{corollary}
(In particular, $\QQ$ is a flat $\ZZ$-module.)
%
%
%
\subsection{Localization Commutes with Everything II}
Let $S$ be a multiplicative set of $A$.
\begin{proposition}[Localization Commutes with Sum
and Intersection]
    If $M$ is an $A$-module and $N$ and $P$
    are submodules of $M$, then
    \begin{alphenum}
    \item $S^{-1}(M/N) = S^{-1}M/S^{-1}N$.
    \item $S^{-1}(N+P) = S^{-1}N + S^{-1}P$.
    \item $S^{-1}(N\cap P) = S^{-1}N\cap S^{-1}P$.
    \end{alphenum}
\end{proposition}
\begin{proposition}[Localization Commutes with Tensor Product and Hom]
    \label{prop:localize_tensor_hom}
    Let $M$ and $N$ be $A$-modules, then
    \begin{alphenum}
    \item 
    $S^{-1}(M\otimes_A N) = S^{-1}M \otimes_{S^{-1}A} S^{-1}N.$
    \item $S^{-1}\Hom_A(M,N)
    = \Hom_{S^{-1}A}(S^{-1}M, S^{-1}N)$.
    \end{alphenum}
\end{proposition}
\begin{exercise}
Prove both of the above theorems.
\end{exercise}
%begin{proof}
 %    -----Proof of (a)-------
 %  Same kind of argument.
 %    Show that one can define maps
 %    \begin{align*}
	% M\otimes N &\to S^{-1}M\otimes_{S^{-1}A}S^{-1}N \\
	% m\otimes n &\mapsto \frac m1 \otimes \frac n1,
 %    \end{align*}
 %    and the inverse
 %    \begin{align*}
	% S^{-1}M\otimes_{S^{-1}A} S^{-1}N
	% &\to S^{-1}(M\otimes_A N) \\
	% \frac ms \otimes \frac nt &\mapsto \frac{m\otimes n}{st}.
	% \qedhere
 %    \end{align*}
%\end{proof}
\section{Local Properties of a Module}
We will now discuss local properties,
which is how properties of localization
transfers to the global ring. The first one is
being the zero module or not.
\begin{proposition}[Being zero module is stalk-local]
    \label{prop:zero_local}
    Let $M$ be an $A$-module. The following are equivalent.
    \begin{alphenum}
	\item $M=0$
	\item $M_{\mathfrak p} = 0$
	    for all prime ideal $\mathfrak p$.
	\item $M_{\mathfrak m}=0$
	    for all maximal ideal $\mathfrak m$.
    \end{alphenum}
\end{proposition}
\begin{proof}
    Clearly (a) $\Rightarrow$ (b) 
    and (b) $\Rightarrow$ (c).
    so we need to check that (c) $\Rightarrow$ (a).

    If $M_{\mathfrak m}=0$ for all maximal ideal $\mathfrak m$,
    then for each $y\in M$, we have $\frac y1=0$
    in $M_{\mathfrak m}$.
    Thus, there exists $a\notin\mathfrak m$
    such that $ay=0\in M$.
    However, this means that $\Ann(y)\not\subset\mathfrak m$
    for each $\mathfrak m$, so by Krull's theorem
    \Cref{thm:krull},
    $\Ann(y) = A$, so $y=0$.
\end{proof}
We say that $M_{\mathfrak p}$ is a \vocab{stalk}
of $M$ at point $[\pp]\in\Spec A$.
Thus, the previous proposition says that
``$M=0$ is a \vocab{stalk-local property}''.
We now discuss more stalk-local properties.
\begin{corollary}[Injectivity 
    and surjectivity is stalk-local]
    \label{cor:injective_surjective_local}
    Let $\phi : M\to N$. Then, the following are equivalent.
    \begin{alphenum}
	\item $\phi : M\to N$ is injective.
	\item $\phi_{\mathfrak p} : M_\pp\to N_\pp$
	    is injective for all prime ideal $\mathfrak p$
	\item $\phi_{\mathfrak m}
    : M_\mm \to N_\mm$
	    is injective for all maxmial ideal $\mathfrak m$.
    \end{alphenum}
    The same statement holds when injective is replaced by
    surjective and isomorphism.
\end{corollary}
\begin{proof}
    Let $K=\Ker\phi$. Then, by exactness of localization
    \Cref{thm:localization_exact}, the exact sequence
    $$0\longto K\longto M \stackrel\phi\longto N$$
    gives the exact sequence
    $$0\longto K_\pp\longto M_\pp\stackrel{\phi_\pp}\longto N_\pp,$$
    so $K_\pp = \Ker\phi_\pp$ for all prime ideal $\pp$.
    Therefore, (b) is equivalent to $K_\pp=0$,
    while (c) is equivalent to $K_\mm=0$.
    Apply \Cref{prop:zero_local} on $K$.

    To prove a similar statement for surjectivity,
    we apply \Cref{prop:zero_local} on $\Coker\phi$ instead.
\end{proof}
The above proposition gives a good way to check 
whether $m_1,\dots,m_n$ generates $M$ or not.
Let $\phi : A^{\oplus n}\to M$ sending the $i$-th
basis element to $m_i$.
With knowledge of localization, we have that
\begin{align*}
    m_1,\dots,m_n\text{ generates } M
    &\iff \phi\text{ is surjective} \\
    &\iff \phi_{\mathfrak m} \text{ is surjective for all }
    \mathfrak m \\
    &\iff \frac{m_1}{1},\dots,\frac{m_n}1\text{ generate }
    M_{\mathfrak m} \\
    &\iff \overline{m_1},\dots,\overline{m_n}
    \text{ generate } M_{\mathfrak m} / \mathfrak m M_{\mathfrak m},
\end{align*}
where the last step follows from Nakayama's lemma
\Cref{cor:nakayama_local_ring}.
This is a convenient procedure to check
if something generates the module
because the last thing is a vector space.


Now, let's discuss some more local properties.

\begin{corollary}[Reduced is stalk-local]
    $A$ is reduced if and only if $A_{\mathfrak m}$ is reduced for
    all maximal ideals $\mathfrak m$.
\end{corollary}
\begin{proof}
    Apply \Cref{prop:zero_local} to $\Nil(A)$,
    noting that localization commutes with nilradical
    \Cref{prop:localization_nilradical}
\end{proof}

\begin{proposition}[Flatness is stalk-local]
    \label{prop:local_flat}
    For any module $M$, $M$ is flat if and only if $M_{\mathfrak m}$ is flat for all maximal ideal
    $\mathfrak m$.
\end{proposition}
\begin{proof}
    The left-to-right direction is easy.
    We prove that the last implies the first.
    Given an injective map $L\to N$,
    then $L_{\mathfrak m}\to N_{\mathfrak m}$
    is injective. Thus, we get the injective map
    $$M_{\mathfrak m}\otimes_{A_{\mathfrak m}}
    L_{\mathfrak m} \to M_{\mathfrak m}\otimes_{A_{\mathfrak m}}
    N_{\mathfrak m}.$$
    However, from the lemma above, we have that
    both sides are isomorphic to
    $$(M\otimes_A L)_{\mathfrak m} \to 
    (M\otimes_A N)_{\mathfrak m}.$$
    Then, since injectivity is local property,
    we get an injective map $M\otimes_A L\to M\otimes_A N$.
    All the maps are natural, so we get the same
    map as what we wanted.
\end{proof}
We have seen many stalk-local properties.
However, not all properties are stalk-local.
\begin{example}[Finitely-generated is not stalk-local]
    Being finitely-generated is not a stalk-local property.
    For example, let $A=\ZZ$ and 
    $$M = \tfrac 12 \ZZ + \tfrac 13 \ZZ + \tfrac 15 \ZZ 
    + \tfrac 17\ZZ + \dots\ \subseteq\ \QQ$$
    (the sum runs through primes).
    Then, for any prime $p$, $M_{(p)} = \tfrac 1p \ZZ_{(p)}$.
    Thus, $M_\mm$ is finitely-generated for all 
    maximal ideal $\mm$.
    However, one can check that $M$ is not finitely-generated.
\end{example}
\section{Locally Free Modules}
\label{sec:locally_free}
\subsection{Definition and Examples}
\begin{definition}[Locally free modules]
    Let $A$ be a ring. A module $M$ is \vocab{locally free}
    if for every prime ideal $\pp$,
    the localization $M_\pp$ is a free $A_\pp$-module.

    A module $M$ is locally free of \vocab{rank} $r$
    if for every prime ideal $\pp$,
    $M_\pp\simeq A_\pp^{\oplus r}$.
\end{definition}
Locally free modules are interesting for two reasons.
\begin{itemize}
    \item They are a natural setting to do multilinear algebra
    (e.g., tensor product, wedge product, dual, trace, determinant).
    \item They encompass a much larger class of examples
    than free modules. For example, if $K$ is 
    an algebraic extension of $\mathbb Q$
    and $\calO_K$ is the ring of integer,
    then any ideal of $\calO_K$ is a locally free
    $\calO_K$-module of rank $1$. 
    This will follow from results in \Cref{chap:dvr}.
\end{itemize}
Consequently, linear algebra over locally free modules
is used in number theory to define norm,
traces, discriminants, etc.
For the rest of this section, we introduce how to do
linear algebra over locally free module.

We do one example that we can simply calculate explicitly.
\begin{example}
    Let $A = \ZZ[\sqrt{-5}]$ and $I=(2,1+\sqrt{-5})$.
    We show that $I$ is locally free.
    To do this, we observe the following.
    \begin{itemize}
    \item If $f_1=2$, then $I_{f_1}$ is locally free
    because $2$ is a unit in $A_{f_1}$,
    so $I_{f_1}$ is a unit ideal.
    \item If $f_2=3$, then $I_{f_2}$ is locally free
    because $2\cdot 3 = (1+\sqrt{-5})(1-\sqrt{-5})$,
    and since $3$ is a unit, it follows that 
    $2$ is divisible by $1+\sqrt{-5}$ in $A_{f_2}$,
    so $I_{f_2}$ is generated by $1+\sqrt{-5}$.
    \end{itemize}
    For every prime ideal $\pp$,
    either $2\notin \pp$ or $3\notin\pp$, 
    so $I_\pp\simeq A_\pp$ for all prime ideal $\pp$.
\end{example}
\subsection{Linear Algebra on Locally Free Modules}
Properties of $\Hom$ and tensor that carries
over from vector spaces (or free modules)
do carry over to locally free modules.
An example is given below.
\begin{theorem}[Linear algebra on locally free modules I]
    Let $P$ and $Q$ be finitely-generated 
    locally free $A$-modules.
    \begin{alphenum}
    \item $P\otimes_A Q$ and $\Hom(P,Q)$ are finitely-generated
    locally free modules.
    \notation{$R^\vee$}{dual of locally free module}
    \item If $R^\vee = \Hom(R,A)$ for all locally free module $R$, 
    then $(P^\vee)^\vee \simeq P$.
    \item For any module $M$, $\Hom(P,M)\simeq P^\vee\otimes_A M$.
    \end{alphenum}
\end{theorem}
\begin{proof}
\begin{alphenum}
    \item Note that for every prime ideal $\pp$,
    $(P\otimes_A Q)_\pp \simeq P_\pp\otimes_{A_\pp} Q_\pp$
    (by \Cref{prop:localize_tensor_hom}) is free 
    because it is a tensor product of free modules.
    Similar reason goes for $\Hom(P,Q)$.
    \item We have a canonical map $P \to (P^\vee)^\vee$ given by
    \begin{align*}
        \phi : P &\to (P^\vee)^\vee \\
        m &\mapsto (f \mapsto f(m)).
    \end{align*}
    By \Cref{cor:injective_surjective_local},
    to show that this is an isomorphism,
    it suffices to show that $\phi_\pp : 
    (P_\pp^\vee)^\vee \to P_\pp$ is an isomorphism
    for all prime ideal $\pp$.
    Thus, it suffices to prove that $\phi$
    is an isomorphism for all finitely-generated free modules $P$,
    i.e., $P\simeq A^{\oplus r}$.
    
    To do this, let $e_1,\dots, e_r$ be the basis.
    Then, observe that
    \begin{itemize}
    \item $f_1,\dots,f_r$ is a basis of $P^\vee$,
    where $f_i(e_j)$ is equal to $1$ if $i=j$ and $0$ otherwise.
    \item $g_1,\dots, g_r$ is a basis of $(P^\vee)^\vee$
    where $g_i(f_j)$ is equal to $1$ if $i=j$
    and $0$ otherwise.
    \end{itemize}
    It is then easy to see that $g_i = \phi(e_i)$,
    so $\phi$ is surjective, hence an isomorphism.
    
    \item Again, we have the canonical map in one direction
    \begin{align*}
        \phi : P^\vee \otimes_A M &\to \Hom(P,M) \\
        f \otimes m &\mapsto (n\mapsto mf(n)).
    \end{align*}
    By the same argument as (b),
    it suffices to check that this is an isomorphism
    for the case when $P$ is a finitely-generated free module. 
    One can do this by computing the basis explicitly,
    similar to above. \qedhere
\end{alphenum}
\end{proof}

We can even define determinants on locally free module.
First, let us recall the abstract definition of determinants.
\begin{definition}[Exterior power]
    \notation{$\bigwedge^n M$}{$n$-th exterior power}
    Let $A$ be a ring and $M$ be an $A$-module.
    A bilinear pairing $\psi :M^n \to P$
    is an \vocab{alternating pairing} if $\psi(x_1,\dots,x_n)=0$
    whenever two of $x_i$ and $x_j$ are equal.
    
    The $n$-th \vocab{exterior power} is
    the module $\bigwedge^n M$ with an alternating pairing
    $\phi : M^n \to \bigwedge^n M$ such that
    any alternating pairing $\psi : M^n\to P$
    factor through $\psi$:
    \begin{center}
    \begin{tikzcd}
        M^n \arrow[rr, "\psi"] 
        \arrow[rd, "\phi"] & & P \\
        & \bigwedge^n M \arrow[ru]
    \end{tikzcd}
    \end{center}
\end{definition}
The explicit construction of exterior power is as follows:
$$\wedgeprod^n M = \frac{M^{\otimes n}}{\langle 
\dots \otimes m \otimes \dots\otimes m\otimes \dots= 0
\text{ for all }m\in M
\rangle},$$
and the image of $(m_1,\dots,m_n)$ under the inclusion
is denoted $m_1\wedge \dots \wedge m_n$.

Observe that since $(x+y)\wedge (x+y)=x\wedge x
= y\wedge y =0$, we get by expanding 
$(x+y)\wedge (x+y)$ that
$$(x\wedge x) + (x\wedge y) + (y\wedge x) + (y\wedge y)
= 0 \implies x\wedge y = -(y\wedge x).$$
(i.e., wedges anticommute).

$\bigwedge^n$ is a functor. Given a morphism 
of modules $\phi : M\to N$, we have the morphism
\begin{align*}
    \wedgeprod^n \phi : \wedgeprod^n M
    &\to \wedgeprod^n N \\
    (m_1\wedge\dots\wedge m_n)
    &\mapsto (\phi(m_1)\wedge\dots\wedge \phi(m_n)).
\end{align*}
\begin{definition}[Determinant]
    \notation{$\det P$}{determinant of locally free module}
    If $P$ is a locally free module of rank $r$,
    then the \vocab{determinant} of $P$
    is the top exterior power $\det P := \bigwedge^r P$,
    which will be a locally free module of rank $1$.
\end{definition}
\begin{example}
    Suppose that $P = A^{\oplus 2}$, with basis $e_1$ and $e_2$. 
    Let $\phi : P\to P$ such that $\phi(e_1)
    = ae_1 + be_2$ and $\phi(e_2) = ce_1+de_2$.
    Then, the map $\bigwedge^2\phi$ sends $e_1\wedge e_2$ 
    to 
    \begin{align*}
        (ae_1+be_2)\wedge (ce_1+de_2)
        &= ac(e_1\wedge e_1) + ad(e_1\wedge e_2)
        + bc(e_2\wedge e_1) + bd(e_2\wedge e_2) \\
        &= ad(e_1\wedge e_2) + bc(e_2\wedge e_1)  \\
        &= (ad-bc)(e_1\wedge e_2).
    \end{align*}
    Thus, the map is multiplication by determinant $ad-bc$.
\end{example}
If $P$ is a free $A$-module, then 
any linear map $\phi : P\to P$ will induce
a scalar multiplication on $\det P\to \det P$,
and this scalar is simply the determinant corresponding
to matrix representing $\phi$.
However, if $P$ is only locally free,
then $\det P$ can be any locally free module of rank $1$.

Many properties of determinants in linear algebra carry over
and can be proven simply by reducing to free modules.
\begin{theorem}[Linear algebra on locally free modules II]
Let $A$ be a ring.
\begin{alphenum}
    \item If $P$ and $Q$ are locally free $A$-modules of rank $r$,
    then $\phi : P\to Q$ is an isomorphism if and only if
    $\bigwedge^r \phi : \det P\to \det Q$ is an isomorphism.

    (This is an analogue of theorem of linear
    algebra that a matrix is invertible if and only
    if the determinant is nonzero.)
    \item If $P$ and $Q$ are locally free $A$-modules
    of rank $p$ and $q$, then
    $$\det(P\oplus Q) \simeq (\det P)\otimes (\det Q).$$
\end{alphenum}
\end{theorem}
\begin{proof}
    One can quickly reduce to the case where $P$ and $Q$ are free. (For (b), one needs to provide the obvious
    map in one direction.)
    We leave this as an exercise to the reader.
\end{proof}
\section{Problems}
\begin{problem}
Construct a ring $A\neq 0$ such that $A_{\mathfrak p}$ is a domain for all prime ideals $\mathfrak p$ but $A$ is not a domain.
\end{problem}
\begin{problem}
Give a new proof of \Cref{prop:nilradical} 
by considering localization $A_f$.
\end{problem}
\begin{problem}
Let $A$ be a ring and $\mathfrak p$ be a prime ideal of $A$ that is minimal under inclusion. Prove that every element of
$\mathfrak p$ is a zero divisor.
\end{problem}
\begin{problem}
\begin{enumerate}[label=(\alph*)]
\item Is the maximal ideal of the local ring $(\mathbb C[x,y]/(y^2-x^3-x))_{(x,y)}$ principal?
\item Is the maximal ideal of the local ring $(\mathbb C[x,y]/(y^2-x^3-x^2))_{(x,y)}$ principal?
\end{enumerate}
\noindent In both parts, we are localizing at maximal ideal $(x,y)$.

(Hint: look at $\mathfrak m/\mathfrak m^2$.)
\end{problem}
\begin{problem}
Let $A$ be a domain and $K=\Frac A$.
The \vocab{rank} of an $A$-module $M$ is defined as 
$\operatorname{rank}_A(M) = \dim_K(M\otimes_A K)$.
\begin{alphenum}
\item What is the rank of the $\ZZ$-module 
$\ZZ^3 \oplus \ZZ/12 \oplus \ZZ/4$?
\item If there is an exact sequence of $A$-modules
$$0\longto M_1 \longto M_2 \longto \dots \longto M_n\longto 0,$$
prove that $\sum_{i=1}^n (-1)^i\operatorname{rank}_A(M_i) = 0$.
\end{alphenum}
\end{problem}
\begin{problem}
Let $A$ be a ring, and let $M$ be a \textbf{finitely-presented} module.
\begin{enumerate}[label=(\alph*)]
\item Prove that if $\mathfrak p$ is a prime ideal such that $M_{\mathfrak p}$ is free, then there exists $f\notin \mathfrak p$ such that $M_f$ is free.%
\footnote{Geometric interpretation: free over point $\mathfrak p$ implies free over open neighborhood $D(f)\ni\mathfrak p$.}
\item Deduce that the following are equivalent.
\begin{itemize}
\item $M_{\mathfrak m}$ is free for all maximal ideals $\mathfrak m$.
\item $M_{\mathfrak p}$ is free for all prime ideals $\mathfrak p$.
\item There exists $f_1,\dots,f_n\in A$ generating a unit ideal such that $M_{f_i}$ is free for all $i=1,2,\dots,n$.
\end{itemize}
\end{enumerate}
\end{problem}
\begin{problem}[Gluing]
\label{prob:gluing}
Let $A$ be a ring, and let $f_1,\dots,f_n\in A$ generating a unit ideal. For each $i$, let $g_i\in A_{f_i}$ such that for all $i,j\in{1,2,\dots,n}$, $g_i$ and $g_j$ map to the same element in $A_{f_if_j}$. Prove that there exists an element $g\in A$ such that $g$ maps to $g_i\in A_{f_i}$ for each $i\in{1,2,\dots,n}$.%
\footnote{This is an algebraic analogue of the following topological fact: given an open cover $U_1,\dots,U_n$ of a topological space $X$ and functions $f_i : U_i\to Y$ such that $f_i|_{U_i\cap U_j} = f_j|_{U_i\cap Uj}$ for all $i,j$, there exists a function $f : X\to Y$ such that $f|_{U_i} = f_i$.}

(Hint: use the fact that $f_1^N, f_2^N,\dots,f_n^N$ generate the unit ideal.) 
\end{problem}
\begin{problem}
Prove that locally free modules are projective
(defined in \Cref{prob:projective}).

(Hint: check condition (ii). Use \Cref{prob:gluing}.)
\end{problem}